% !TeX root = main.tex

\chapter{Temas Avanzados}

\begin{quotation}
	\itshape
	``La geometría de la información no termina con las familias exponenciales; se extiende a variedades más generales, conexiones no lineales y estructuras que aún están siendo descubiertas.''
\end{quotation}

En este capítulo exploramos extensiones y generalizaciones de la geometría de la información.  Vamos más allá de las familias exponenciales para estudiar \textbf{variedades estadísticas generales}, \textbf{geometría de información cuántica}, \textbf{aprendizaje de variedades} y \textbf{estructuras emergentes} como la geometría de información alpha-beta y las conexiones de Wasserstein.

\section{Variedades estadísticas generales}

Hasta ahora nos enfocamos en familias exponenciales, donde la geometría es particularmente limpia.  Pero la mayoría de distribuciones no son exponenciales.

\begin{definition}[Variedad estadística general]
	Una \textbf{variedad estadística} es un conjunto de distribuciones de probabilidad $\mathcal{M} = \{p_\theta : \theta \in \Theta\}$ parametrizado suavemente, donde $\Theta$ es un conjunto abierto de $\mathbb{R}^k$.
\end{definition}

\begin{proposition}[Métrica de Fisher general]
	Para una variedad estadística general, la métrica de Fisher $g_{ij}(\theta) = \mathbb{E}[\partial_i \ell \cdot \partial_j \ell]$ siempre está bien definida, independientemente de si la familia es exponencial o no.
\end{proposition}

\begin{example}[Familia t-Student]
	La distribución t-Student con $\nu$ grados de libertad y parámetro de localización $\mu$:
	\[
	p_\mu(x) = \frac{\Gamma((\nu+1)/2)}{\sqrt{\nu\pi}\,\Gamma(\nu/2)} \left(1 + \frac{(x-\mu)^2}{\nu}\right)^{-(\nu+1)/2}.
	\]
	Esta familia \emph{no} es exponencial (no hay estadístico suficiente de dimensión finita), pero la métrica de Fisher está bien definida.
\end{example}

\begin{center}
	\begin{tikzpicture}
		\node[draw, rounded corners, fill=AccentLight!20, minimum width=2.5cm, minimum height=1cm] (exp) at (0,0) {Exponenciales};
		\node[draw, rounded corners, fill=Accent!20, minimum width=2.5cm, minimum height=1cm] (gen) at (4,0) {Generales};
		\node[below=0.3cm] at (exp.south) {\scriptsize Geometría limpia};
		\node[below=0.3cm] at (gen.south) {\scriptsize Geometría rica};
		\draw[->, thick, AccentDark] (exp) -- (gen) node[midway, above] {\scriptsize generalización};
	\end{tikzpicture}
\end{center}

\section{Geometría de información cuántica}

La geometría de la información se extiende naturalmente al caso cuántico, donde las ``distribuciones'' son \textbf{estados cuánticos} (operadores densidad).

\begin{definition}[Estado cuántico]
	Un \textbf{estado cuántico} es un operador densidad $\rho$ que satisface: $\rho \ge 0$, $\operatorname{Tr}(\rho) = 1$.
\end{definition}

\begin{definition}[Métrica de Bures]
	La \textbf{métrica de Bures} es la métrica cuántica natural:
	\[
	d_B(\rho, \sigma) = \arccos\left(\operatorname{Tr}\sqrt{\sqrt{\rho}\,\sigma\sqrt{\rho}}\right).
	\]
	Es la generalización cuántica de la métrica de Fisher--Rao.
\end{definition}

\begin{theorem}[Relación con la métrica de Fisher]
	Para estados cuánticos puros $|\psi\rangle$, la métrica de Bures reduce a la métrica de Fubini--Study, que es la métrica de Fisher--Rao para amplitudes de probabilidad.
\end{theorem}

\section{Aprendizaje de variedades}

\begin{definition}[Aprendizaje de variedades]
	El \textbf{aprendizaje de variedades} (manifold learning) es el problema de reconstruir una variedad de baja dimensión incrustada en un espacio de alta dimensión a partir de datos observados.
\end{definition}

\begin{proposition}[Conexión con la geometría de la información]
	Cuando los datos son muestras de una distribución, la variedad aprendida puede interpretarse como una \textbf{variedad estadística}, y los métodos de aprendizaje de variedades (Isomap, LLE, t-SNE) son aproximaciones de geodésicas en la métrica de Fisher.
\end{proposition}

\section{Geometría de information $\alpha$-$\beta$}

\begin{definition}[$\alpha$-$\beta$-divergencia]
	La \textbf{$\alpha$-$\beta$-divergencia} generaliza la KL y las $f$-divergencias:
	\[
	D^{(\alpha,\beta)}(p \| q) = \frac{4}{1-\alpha^2}\left(1 - \sum_x p(x)^{(1+\alpha)/2} q(x)^{(1-\alpha)/2}\right)
	\]
	para $\alpha \ne \pm 1$, con extensiones por continuidad.
\end{definition}

\begin{theorem}[Propiedades]
	\begin{enumerate}
		\item Para $\alpha = \beta = 0$: $D^{(0,0)} = D_{\mathrm{KL}}$ (divergencia KL).
		\item Para $\alpha = 1$: $D^{(1,\beta)}$ es una $f$-divergencia.
		\item Para $\beta = 0$: $D^{(\alpha,0)}$ es una divergencia de Bregman.
	\end{enumerate}
\end{theorem}

\section{Conexiones de Wasserstein}

\begin{definition}[Distancia de Wasserstein]
	La \textbf{distancia de Wasserstein} de orden 2 entre dos distribuciones $p$ y $q$ en $\mathbb{R}^d$ es:
	\[
	W_2(p, q) = \inf_{\gamma \in \Pi(p,q)} \left(\int \|x - y\|^2 \, d\gamma(x,y)\right)^{1/2}
	\]
	donde $\Pi(p,q)$ es el conjunto de acoplamientos (joint distributions) con marginales $p$ y $q$.
\end{definition}

\begin{proposition}[Ventajas de Wasserstein]
	La distancia de Wasserstein tiene ventajas sobre la KL:
	\begin{itemize}
		\item Es una métrica verdadera (simétrica, satisface triangular).
		\item Está bien definida incluso cuando las distribuciones no tienen soporte común.
		\item Codifica la ``geometría del espacio muestral'' (no solo la probabilidad).
	\end{itemize}
\end{proposition}

\section{Optimal Transport y geometría}

\begin{theorem}[Kantorovich--Rubinstein]
	La distancia de Wasserstein de orden 1 puede expresarse como:
	\[
	W_1(p, q) = \sup_{\|f\|_L \le 1} \left(\mathbb{E}_p[f(x)] - \mathbb{E}_q[f(x)]\right)
	\]
	donde el supremo es sobre funciones 1-Lipschitz.
\end{theorem}

\begin{example}[Transporte óptimo en imágenes]
	En procesamiento de imágenes, el transporte óptimo se usa para comparar histogramas de color, transferir estilos y alinear registros médicos.
\end{example}

\section{Geometría de información y aprendizaje profundo}

\subsection{El espacio de parámetros de una red neuronal}

\begin{proposition}[Geometría del loss landscape]
	El loss landscape de una red neuronal puede interpretarse como una \textbf{variedad estadística} con métrica de Fisher inducida por los datos.  Los mínimos locales corresponden a ``valles'' en esta variedad.
\end{proposition}

\subsection{Batch normalization y geometría}

\begin{theorem}[Batch normalization como proyección]
	La batch normalization es equivalente a una proyección $e$-exponencial de las activaciones sobre la familia de distribuciones normalizadas.
\end{theorem}

\section{Código Python: transporte óptimo}

\begin{lstlisting}[language=Python, style=PythonStyle]
import numpy as np
from scipy.optimize import linprog

def earth_movers_distance_1d(p, q, x):
    """Calcula la distancia de Wasserstein 1D entre distribuciones discretas."""
    # Funcion de distribucion acumulada
    cdf_p = np.cumsum(p)
    cdf_q = np.cumsum(q)
    
    # W1 = integral de |F_p - F_q|
    w1 = np.trapz(np.abs(cdf_p - cdf_q), x)
    return w1

# Ejemplo
x = np.array([0, 1, 2, 3, 4])
p = np.array([0.1, 0.2, 0.4, 0.2, 0.1])
q = np.array([0.2, 0.1, 0.1, 0.4, 0.2])

w1 = earth_movers_distance_1d(p, q, x)

# Verificacion con KL
from scipy.stats import entropy
p_smooth = p + 1e-10
q_smooth = q + 1e-10
kl = entropy(p_smooth, q_smooth)
# Result: Wasserstein-1 \approx  0.6, KL \approx  0.36
\end{lstlisting}

\begin{lstlisting}[language=Python, style=PythonStyle]
# Geodesicas en el espacio de Gaussianas
def gaussian_geodesic(mu0, sigma0, mu1, sigma1, t):
    """Geodesica entre dos Gaussianas en coordenadas de Fisher."""
    # Coordenadas naturales
    eta0 = np.array([mu0/sigma0**2, -1/(2*sigma0**2)])
    eta1 = np.array([mu1/sigma1**2, -1/(2*sigma1**2)])
    
    # Interpolacion en eta (geodesica exponencial)
    eta_t = (1 - t) * eta0 + t * eta1
    
    # Conversion a mu, sigma
    sigma_t = np.sqrt(-1 / (2 * eta_t[1]))
    mu_t = eta_t[0] * sigma_t**2
    
    return mu_t, sigma_t

# Ejemplo
mu0, sigma0 = 0, 1
mu1, sigma1 = 3, 2

for t in [0, 0.25, 0.5, 0.75, 1.0]:
    mu_t, sigma_t = gaussian_geodesic(mu0, sigma0, mu1, sigma1, t)
# Result: t=0 => (0,1), t=0.5 => (1.5, 1.414), t=1 => (3,2)
\end{lstlisting}

\section{Ejercicios}

\subsection*{Variedades generales}
\begin{enumerate}
	\item[$\star$] Defina una variedad estadística y explique por qué las familias exponenciales son un caso especial.
	\item[$\star\star$] Para la familia t-Student con $\nu = 5$ y parámetro de localización $\mu$, calcule la métrica de Fisher.
	\item[$\star$] ¿Por qué la familia t-Student no es exponencial?
	\item[$\star\star\star$] Demuestre que la métrica de Fisher está bien definida para cualquier familia paramétrica regular.
	\item[$\star\star$] Compare la métrica de Fisher de la Normal con la de la t-Student para $\nu = 5$.  ¿En qué regiones difieren más?
	\item[$\star\star\star$] Para la familia de Cauchy, muestre que la métrica de Fisher es constante e independiente del parámetro de localización.
	\item[$\star\star\star$] Discuta por qué la familia uniforme $U(0, \theta)$ no tiene una métrica de Fisher ``natural'' en el sentido de Chentsov.
\end{enumerate}

\subsection*{Geometría cuántica}
\begin{enumerate}[resume]
	\item[$\star$] Defina un estado cuántico (operador densidad) y la métrica de Bures.
	\item[$\star\star$] Para un qubit con estado $|\psi\rangle = \cos(\theta/2)|0\rangle + \sin(\theta/2)e^{i\phi}|1\rangle$, calcule la métrica de Bures.
	\item[$\star\star\star$] Muestre que la métrica de Bures reduce a la métrica de Fubini--Study para estados puros.
	\item[$\star$] ¿Cuál es la interpretación física de la distancia de Bures?
	\item[$\star\star\star$] Para un sistema de dos qubits, describa por qué la geometría del espacio de estados entrelazados es fundamentalmente diferente a la de estados separables.
	\item[$\star\star$] Compare la métrica de Bures con la distancia trace: ¿cuándo coinciden y cuándo difieren?
\end{enumerate}

\subsection*{Transporte óptimo}
\begin{enumerate}[resume]
	\item[$\star$] Defina la distancia de Wasserstein $W_1$ y explique qué mide geométricamente.
	\item[$\star\star$] Calcule la distancia de Wasserstein $W_1$ entre $\mathrm{Bern}(0.3)$ y $\mathrm{Bern}(0.7)$.
	\item[$\star\star$] Para distribuciones discretas $p = (0.5, 0.5)$ y $q = (0.2, 0.8)$ en $\{0, 1\}$, encuentre el transporte óptimo.
	\item[$\star\star\star$] Muestre que $W_1$ es una métrica verdadera (simétrica, definitiva positiva, triangular).
	\item[$\star\star\star\star$] (Reto) Demuestre el teorema de Kantorovich--Rubinstein para el caso unidimensional.
	\item[$\star\star\star$] Para $\mathrm{Uniforme}(0,1)$ y $\mathrm{Uniforme}(0,2)$, calcule $W_2^2$ y compare con la divergencia KL.
	\item[$\star\star$] Explique por qué la distancia de Wasserstein es preferible a la KL para comparar histogramas de imágenes.
	\item[$\star\star\star$] Para distribuciones Gaussianas univariadas, muestre que $W_2^2(\mathcal{N}(\mu_1, \sigma_1^2), \mathcal{N}(\mu_2, \sigma_2^2)) = (\mu_1 - \mu_2)^2 + (\sigma_1 - \sigma_2)^2$.
\end{enumerate}

\subsection*{Temas emergentes}
\begin{enumerate}[resume]
	\item[$\star$] Explique la conexión entre batch normalization y proyección exponencial.
	\item[$\star$] Para una red neuronal con activación softmax, ¿cuál es la métrica de Fisher inducida?
	\item[$\star$] ¿Cómo se relaciona la geometría de Wasserstein con la geometría de Fisher?
	\item[$\star\star\star$] Para un modelo de lenguaje con softmax, muestre que la perplejidad está relacionada con la divergencia KL entre la distribución de datos y la distribución del modelo.
	\item[$\star\star\star$] Discuta por qué el gradiente natural es computacionalmente costoso para redes neuronales profundas y cómo las aproximaciones (K-FAC) mitigan este problema.
	\item[$\star\star\star\star$] (Reto) Proponga una extensión de la divergencia KL para distribuciones con soportes diferentes usando transporte óptimo.
\end{enumerate}

\subsection*{Qué gana con GI}
\begin{quote}
	\itshape
	Sin GI, la geometría de distribuciones se trata como una curiosidad matemática.  Con GI, se descubre que la misma estructura---métrica de Fisher, dualidad de Legendre, proyecciones KL---aparece en dominios aparentemente dispares: credibilidad actuarial, optimización de redes neuronales, inferencia bayesiana aproximada y transporte óptimo.  La geometría no es una herramienta más en el cinturón; es el hilo conductor que explica por qué ciertos métodos funcionan y cómo mejorarlos de manera sistemática, no por prueba y error.
\end{quote}
